\section{Conclusión}

\begin{itemize}
    \item La implementación de Odoo en Industria Textil Atlas E.I.R.L. ha permitido erradicar de forma definitiva los procesos manuales y la fragmentación de la información comercial y logística. La centralización del catálogo maestro de 22 artículos deportivos bajo códigos SKU asegura que almacén, mostrador y compras compartan datos consistentes en la nube en tiempo real.
    \item Se ha logrado agilizar significativamente la venta y atención al cliente en el mostrador físico de la tienda de San Camilo. Gracias a la terminal del Punto de Venta (POS) y las cotizaciones mayoristas rápidas en Ventas, el tiempo de preparación de pedidos y facturación se redujo de 15--20 minutos por cliente a menos de 5 minutos, eliminando las esperas y los errores de cálculo.
    \item La digitalización del flujo de compras en base a reglas de stock mínimo automático ha dotado a la gerencia general de Víctor Huayllaro de una herramienta financiera robusta. Esto permite planificar el aprovisionamiento con los talleres confeccionistas externos basándose en datos de demanda reales, resguardando la liquidez y controlando los márgenes comerciales de forma estricta.
    \item El programa de capacitación práctica y simulación previa en base de datos de pruebas (Sandbox) logró superar con éxito el bajo nivel de madurez digital inicial del personal de mostrador (35.34\%), logrando una tasa de adopción activa superior al 90\% en las tareas operativas cotidianas del negocio.
    \item Finalmente, el proyecto ha sentado las bases tecnológicas e informativas necesarias para la escalabilidad del negocio en Arequipa. Con una infraestructura digital integrada y un control riguroso de existencias, la empresa se encuentra en una posición inmejorable para abordar en una fase posterior la unificación definitiva de la facturación electrónica SUNAT directamente dentro del sistema Odoo, eliminando la necesidad de operar un doble sistema paralelo.
\end{itemize}
